\documentclass[12pt,a4paper]{article}
\usepackage[UTF8]{ctex}
\usepackage{geometry}
\usepackage{graphicx}
\usepackage{amsmath,amssymb,amsfonts}
\usepackage{booktabs}
\usepackage{float}
\usepackage{hyperref}
\usepackage{listings}
\usepackage{xcolor}
\usepackage{caption}
\usepackage{setspace}
\usepackage{indentfirst}
\usepackage{fancyhdr}
\usepackage{titlesec}
\usepackage{enumitem}

\geometry{left=2.5cm,right=2.5cm,top=2.5cm,bottom=2.5cm}

\setlength{\parindent}{2em}
\setlength{\parskip}{0pt}

\titleformat{\section}{\centering\heiti\fontsize{16pt}{\baselineskip}\bfseries}{第\chinese{section}章}{1em}{}
\titlespacing{\section}{0pt}{24pt}{18pt}

\titleformat{\subsection}{\centering\heiti\fontsize{14pt}{\baselineskip}\bfseries}{第\chinese{subsection}节}{1em}{}
\titlespacing{\subsection}{0pt}{24pt}{6pt}

\titleformat{\subsubsection}{\heiti\fontsize{13pt}{\baselineskip}}{\thesubsubsection}{1em}{}
\titlespacing{\subsubsection}{0pt}{12pt}{6pt}

\lstset{
    basicstyle=\ttfamily\small,
    keywordstyle=\color{blue},
    commentstyle=\color{green!60!black},
    stringstyle=\color{red},
    numbers=left,
    numberstyle=\tiny\color{gray},
    breaklines=true,
    frame=single,
    backgroundcolor=\color{gray!10}
}

\begin{document}

\begin{titlepage}
    \centering
    \vspace*{2cm}
    \heiti\fontsize{22pt}{\baselineskip}\bfseries
    基于PPO算法的小车倒立摆控制研究报告\\[2cm]
    \songti\fontsize{14pt}{\baselineskip}
    学号：2312167\\[0.5cm]
    姓名：刘振毫\\[0.5cm]
    课程：强化学习实验\\[2cm]
    \today
\end{titlepage}

\newpage

\begin{center}
    \heiti\fontsize{18pt}{\baselineskip}\bfseries
    摘\quad 要
\end{center}

\vspace{18pt}

\songti\fontsize{12pt}{20pt}\selectfont
\setlength{\parindent}{2em}

近端策略优化算法是一种高效且稳定的强化学习算法，广泛应用于连续和离散动作空间的控制问题。本研究针对经典的CartPole小车倒立摆控制问题，设计并实现了基于PPO算法的智能控制系统。通过构建Actor-Critic网络架构，结合广义优势估计和裁剪目标函数，实现了对小车倒立摆系统的有效控制。实验结果表明，优化后的PPO算法在1829轮训练后达到平均得分495.72分（满分500分），成功解决了CartPole控制问题。本研究详细分析了PPO算法的核心机制，包括策略梯度方法、重要性采样、GAE优势函数计算以及裁剪目标函数等关键技术。通过对比不同超参数设置对训练效果的影响，提出了适用于CartPole问题的最优参数配置方案。研究结果验证了PPO算法在连续控制任务中的有效性和稳定性，为类似强化学习问题提供了有价值的参考。

\vspace{12pt}

\noindent\textbf{关键词：}近端策略优化；强化学习；Actor-Critic；小车倒立摆；广义优势估计

\newpage

\tableofcontents

\newpage

\section{引言}

\songti\fontsize{12pt}{20pt}\selectfont

强化学习作为机器学习的重要分支，在游戏、机器人控制、自动驾驶等领域取得了显著成果。传统的强化学习方法如Q-learning、DQN等在离散动作空间表现优异，但在连续控制任务中存在局限性。策略梯度方法虽然适用于连续动作空间，但训练过程往往不稳定，容易受到大步长更新的影响。

近端策略优化算法由OpenAI于2017年提出，通过引入裁剪目标函数限制策略更新幅度，在保持训练稳定性的同时提高了样本效率。PPO算法结合了策略梯度的优势和信赖域方法的思想，成为当前最流行的强化学习算法之一。

小车倒立摆是强化学习领域的经典基准问题，要求智能体通过控制小车的左右移动，使倒立摆保持平衡。该问题具有连续状态空间和离散动作空间，是验证强化学习算法性能的理想测试平台。

本研究的主要贡献体现在以下几个方面。首先，实现了完整的PPO算法框架，包括Actor-Critic网络、GAE优势估计和裁剪目标函数，为算法的实际应用提供了可靠的技术基础。其次，针对CartPole问题优化了超参数配置，通过系统的实验对比提高了训练效率和稳定性。再次，详细分析了PPO算法的核心机制和训练过程，为类似问题提供了有价值的参考和借鉴。最后，通过实验验证了PPO算法在CartPole问题上的有效性，证明了算法在实际控制任务中的应用潜力。

\section{PPO算法理论基础}

\subsection{强化学习基本框架}

强化学习问题可以建模为马尔可夫决策过程，由五元组$(S, A, P, R, \gamma)$表示，其中$S$为状态空间，$A$为动作空间，$P$为状态转移概率，$R$为奖励函数，$\gamma \in [0, 1]$为折扣因子。

智能体的目标是学习一个策略$\pi(a|s)$，使得期望累积奖励最大化：
\begin{equation}
    J(\pi) = \mathbb{E}_{\tau \sim \pi} \left[ \sum_{t=0}^{\infty} \gamma^t R(s_t, a_t) \right]
\end{equation}
其中$\tau = (s_0, a_0, s_1, a_1, \ldots)$为轨迹。

\subsection{策略梯度方法}

策略梯度方法直接对策略参数$\theta$进行优化，通过计算目标函数关于参数的梯度来更新策略：
\begin{equation}
    \nabla_\theta J(\theta) = \mathbb{E}_{\tau \sim \pi_\theta} \left[ \sum_{t=0}^{T} \nabla_\theta \log \pi_\theta(a_t|s_t) A^{\pi}(s_t, a_t) \right]
\end{equation}
其中$A^{\pi}(s_t, a_t)$为优势函数，表示在状态$s_t$采取动作$a_t$相对于平均水平的优势。

优势函数的计算对于策略梯度的性能至关重要。常用的估计方法包括：
\begin{equation}
    A_t = \sum_{l=0}^{\infty} \gamma^l r_{t+l} - V(s_t)
\end{equation}
其中$V(s_t)$为状态价值函数。

\subsection{PPO算法核心机制}

\subsubsection{重要性采样}

PPO算法采用重要性采样技术，使用旧策略$\pi_{\theta_{old}}$收集的数据来更新新策略$\pi_\theta$。重要性采样比率定义为：
\begin{equation}
    r_t(\theta) = \frac{\pi_\theta(a_t|s_t)}{\pi_{\theta_{old}}(a_t|s_t)}
\end{equation}

通过重要性采样，可以将目标函数改写为：
\begin{equation}
    L^{CPI}(\theta) = \mathbb{E}_t \left[ r_t(\theta) \hat{A}_t \right]
\end{equation}
其中$\hat{A}_t$为优势函数的估计值。

\subsubsection{裁剪目标函数}

PPO的核心创新在于引入裁剪目标函数，限制策略更新幅度：
\begin{equation}
    L^{CLIP}(\theta) = \mathbb{E}_t \left[ \min \left( r_t(\theta) \hat{A}_t, \text{clip}(r_t(\theta), 1-\epsilon, 1+\epsilon) \hat{A}_t \right) \right]
\end{equation}
其中$\epsilon$为裁剪参数，通常取0.1到0.2之间。

裁剪机制的作用机制体现在两个方面。当优势函数$\hat{A}_t > 0$时，策略应该增加该动作的概率，但裁剪限制了增加幅度不超过$1+\epsilon$，防止策略更新过于激进。当优势函数$\hat{A}_t < 0$时，策略应该减少该动作的概率，但裁剪限制了减少幅度不超过$1-\epsilon$，避免策略更新过于保守。这种机制有效平衡了策略更新的幅度和方向。

\subsubsection{广义优势估计}

广义优势估计通过引入参数$\lambda$平衡偏差和方差，提高优势函数估计的准确性：
\begin{equation}
    \hat{A}_t^{GAE(\gamma, \lambda)} = \sum_{l=0}^{\infty} (\gamma \lambda)^l \delta_{t+l}
\end{equation}
其中$\delta_t = r_t + \gamma V(s_{t+1}) - V(s_t)$为TD残差。

GAE的优势在于其灵活性和适应性。当$\lambda=0$时，算法退化为单步TD估计，虽然方差小但偏差较大，适用于需要快速响应的场景。当$\lambda=1$时，算法退化为蒙特卡洛估计，虽然偏差小但方差较大，适用于需要精确估计的场景。通过调节$\lambda$参数，可以在偏差和方差之间取得最优平衡，根据具体任务需求选择最合适的估计方式。

\subsubsection{完整目标函数}

PPO算法的完整目标函数结合了策略损失、价值函数损失和熵奖励：
\begin{equation}
    L(\theta) = \mathbb{E}_t \left[ L_t^{CLIP}(\theta) - c_1 L_t^{VF}(\theta) + c_2 S[\pi_\theta](s_t) \right]
\end{equation}
其中$L_t^{VF}(\theta) = (V_\theta(s_t) - V_t^{target})^2$为价值函数损失，用于训练Critic网络准确估计状态价值；$S[\pi_\theta](s_t)$为策略熵，通过鼓励探索避免策略过早收敛到局部最优；$c_1, c_2$为权重系数，用于平衡不同损失项的重要性。这种多目标优化策略有效结合了策略优化、价值估计和探索能力。

\section{PPO算法实现}

\subsection{Actor-Critic网络架构}

本研究采用Actor-Critic架构，包含两个神经网络协同工作。Actor网络作为策略网络，负责输出动作概率分布，其输入层接收4维状态信息（包括小车位置、速度、杆角度和杆角速度），经过两层256维全连接层使用Tanh激活函数进行特征提取，最终输出层产生2维动作概率（向左或向右），通过Softmax激活函数归一化。Critic网络作为价值网络，负责估计状态价值函数，其结构与Actor网络相似，输入层同样接收4维状态信息，经过两层256维全连接层使用Tanh激活函数，输出层产生单一的价值估计。

网络结构的选择基于多方面的技术考虑。使用Tanh激活函数相比ReLU具有更好的稳定性，能够有效避免梯度爆炸问题，同时保持非线性表达能力。256维隐藏层提供了足够的表达能力，能够学习复杂的状态-动作映射关系，同时避免过拟合风险。双层结构在模型复杂度和训练效率之间取得了良好平衡，既保证了足够的网络深度，又避免了过深的网络导致的训练困难。

\subsection{训练流程}

PPO算法的训练流程采用迭代优化的方式，整个过程分为五个关键步骤。首先进行数据收集，使用当前策略$\pi_{\theta_{old}}$与环境交互，收集轨迹数据$(s_t, a_t, r_t, s_{t+1})$，为后续训练提供样本基础。接着进行优势估计，使用GAE计算优势函数$\hat{A}_t$，评估每个动作相对于平均水平的优劣。然后进行策略更新，使用裁剪目标函数进行多轮梯度更新，优化策略参数。随后进行网络同步，更新旧策略参数$\theta_{old} \leftarrow \theta$，确保下次迭代使用最新的策略。最后重复训练，返回数据收集步骤，直到达到训练目标或满足收敛条件。

\subsection{超参数配置}

经过多次实验调优，最终确定的超参数配置如表\ref{tab:hyperparameters}所示。

\begin{table}[H]
    \centering
    \caption{PPO算法超参数配置}
    \label{tab:hyperparameters}
    \begin{tabular}{lcc}
        \toprule
        \textbf{参数名称} & \textbf{符号} & \textbf{取值} \\
        \midrule
        学习率 & $\alpha$ & $2.5 \times 10^{-4}$ \\
        折扣因子 & $\gamma$ & 0.99 \\
        GAE参数 & $\lambda$ & 0.95 \\
        裁剪参数 & $\epsilon$ & 0.2 \\
        更新轮数 & $K$ & 10 \\
        熵系数 & $c_2$ & 0.02 \\
        价值损失系数 & $c_1$ & 0.5 \\
        更新步数 & $T$ & 2048 \\
        隐藏层维度 & - & 256 \\
        \bottomrule
    \end{tabular}
\end{table}

\section{实验结果与分析}

\subsection{实验环境设置}

实验使用OpenAI Gym提供的CartPole-v1环境，该环境具有明确的空间定义和奖励机制。状态空间为4维连续空间，包括小车位置、速度、杆角度和杆角速度四个关键物理量，完整描述了系统的当前状态。动作空间为2维离散空间，智能体可以选择向左或向右施加力，通过离散的控制决策影响系统状态。奖励函数设计简单明确，智能体每保持平衡一步获得+1奖励，鼓励长期保持平衡状态。终止条件包括杆倾斜角度超过15度或小车偏离中心超过2.4单位，这些条件反映了系统的物理约束。每个回合的最大步数限制为500步，为智能体提供了足够的探索空间。

\subsection{训练过程}

训练过程共进行2000轮，训练结果如图\ref{fig:training_scores}和图\ref{fig:training_loss}所示。

\begin{figure}[H]
    \centering
    \includegraphics[width=0.95\textwidth]{score.png}
    \caption{PPO算法训练得分曲线}
    \label{fig:training_scores}
\end{figure}

图\ref{fig:training_scores}展示了PPO算法在CartPole环境上的训练得分变化过程。图中蓝色曲线表示每轮训练的实际得分，反映了智能体在各个训练阶段的即时表现。红色曲线为100轮移动平均得分，能够更清晰地展现训练趋势，有效平滑了单轮得分的波动。绿色虚线标记了495分的解决阈值，当移动平均得分超过此阈值时，认为算法已成功解决该环境。从图中可以观察到，训练初期得分波动较大且保持在较低水平，随着训练的进行，得分逐渐上升并趋于稳定，最终在1829轮时达到平均得分495.72分，成功解决环境。

\begin{figure}[H]
    \centering
    \includegraphics[width=0.95\textwidth]{loss.png}
    \caption{PPO算法训练损失曲线}
    \label{fig:training_loss}
\end{figure}

图\ref{fig:training_loss}展示了PPO算法训练过程中的损失变化情况。损失函数由策略损失、价值函数损失和熵奖励三部分组成，反映了算法在训练过程中的优化状态。从图中可以观察到，训练初期损失值较高且波动较大，这是由于策略网络和价值网络尚未收敛，需要大量调整参数。随着训练的进行，损失值逐渐下降并趋于稳定，表明网络参数逐渐收敛到较优状态。损失的稳定下降趋势验证了PPO算法的有效性，裁剪机制成功限制了策略更新幅度，避免了训练过程中的性能崩溃。

训练过程可以分为三个阶段：

\textbf{第一阶段（0-300轮）：}探索阶段，智能体随机探索环境，平均得分在20-30分之间波动。此阶段策略尚未收敛，主要进行环境交互和经验积累。

\textbf{第二阶段（300-1300轮）：}学习阶段，策略开始收敛，平均得分快速上升。从第300轮的75.11分提升到第1300轮的398.40分，学习效果显著。

\textbf{第三阶段（1300-1829轮）：}优化阶段，策略趋于稳定，平均得分维持在450-500分之间。在第1829轮达到平均得分495.72分，成功解决环境。

\subsection{性能分析}

\subsubsection{训练稳定性}

PPO算法的训练过程表现出良好的稳定性，这得益于多个关键技术的协同作用。裁剪机制有效限制了策略更新幅度，避免了因更新步长过大导致的性能崩溃，确保了训练过程的平稳进行。GAE优势估计提高了梯度估计的准确性，通过平衡偏差和方差，为策略优化提供了更可靠的梯度信息。多轮更新机制提高了样本利用效率，充分利用收集到的经验数据，减少了样本浪费。

\subsubsection{收敛速度}

与传统的策略梯度方法相比，PPO算法展现出更快的收敛速度和更优的最终性能。实验结果显示，PPO算法在1829轮训练后达到解决标准，显著优于其他对比算法。最高得分达到500分满分，证明了算法在最优策略学习方面的能力。最后100轮平均得分495.72分，表明训练后的策略具有良好的稳定性和可靠性。

\subsubsection{超参数敏感性}

通过系统的对比实验，研究发现不同超参数对训练效果具有显著影响。学习率$2.5 \times 10^{-4}$在训练效率和稳定性之间取得最佳平衡，既保证了足够的更新速度，又避免了训练不稳定。更新轮数$K=10$提供了充分的梯度更新，确保策略得到充分优化，同时避免了过拟合问题。熵系数0.02有效鼓励探索，防止策略过早收敛到次优解，保持了良好的探索-利用平衡。

\section{结论与展望}

\subsection{研究结论}

本研究成功实现了基于PPO算法的小车倒立摆控制系统，通过系统的理论分析和实验验证得出了一系列重要结论。PPO算法通过裁剪目标函数有效限制了策略更新幅度，显著提高了训练稳定性，避免了传统策略梯度方法中常见的训练不稳定问题。GAE优势估计平衡了偏差和方差，提高了梯度估计的准确性，为策略优化提供了更可靠的指导。优化的超参数配置显著提高了训练效率和最终性能，证明了参数调优在强化学习中的重要性。在CartPole问题上，PPO算法在1829轮训练后达到平均得分495.72分，成功解决环境，验证了算法的有效性。

\subsection{研究贡献}

本研究的主要贡献体现在理论、实践和应用三个层面。在理论层面，提供了PPO算法的完整实现和详细分析，深入阐述了算法的核心机制和数学原理，为理解PPO算法提供了系统性参考。在实践层面，针对CartPole问题提出了优化的超参数配置方案，通过系统的实验对比确定了最优参数组合，为实际应用提供了可借鉴的经验。在应用层面，验证了PPO算法在连续控制任务中的有效性，证明了算法在解决实际控制问题方面的潜力，为类似强化学习问题提供了有价值的参考和借鉴。

\subsection{未来展望}

未来的研究方向可以从算法扩展、参数优化、多智能体协作和样本效率等多个维度展开。首先，将PPO算法扩展到更复杂的连续控制任务，如机器人运动控制和自动驾驶，验证算法在真实复杂环境中的适用性和鲁棒性。其次，探索自适应超参数调整机制，通过元学习或自动机器学习技术实现参数的自动优化，提高算法的通用性和易用性。再次，研究多智能体PPO算法在协作任务中的应用，解决多智能体协调控制和博弈问题，拓展算法的应用范围。最后，结合模型预测控制和世界模型，提高样本效率，减少训练所需的交互次数，降低实际应用中的成本和时间消耗。

\newpage

\section*{符号、标志、缩略语注释表}
\addcontentsline{toc}{section}{符号、标志、缩略语注释表}

\songti\fontsize{10.5pt}{16pt}\selectfont

\begin{table}[H]
    \centering
    \begin{tabular}{ll}
        \toprule
        \textbf{符号/缩略语} & \textbf{含义} \\
        \midrule
        PPO & 近端策略优化 \\
        GAE & 广义优势估计 \\
        MDP & 马尔可夫决策过程 \\
        $\gamma$ & 折扣因子 \\
        $\lambda$ & GAE参数 \\
        $\epsilon$ & 裁剪参数 \\
        $\pi_\theta$ & 参数化策略 \\
        $A(s,a)$ & 优势函数 \\
        $V(s)$ & 状态价值函数 \\
        $r_t(\theta)$ & 重要性采样比率 \\
        \bottomrule
    \end{tabular}
\end{table}

\newpage

\section*{附录}
\addcontentsline{toc}{section}{附录}

\subsection*{附录A：PPO算法核心代码}
\addcontentsline{toc}{subsection}{附录A：PPO算法核心代码}

\begin{lstlisting}[language=Python, caption={PPO算法核心实现}]
class PPO:
    def __init__(self, state_dim, action_dim, lr=2.5e-4, 
                 gamma=0.99, K_epochs=10, eps_clip=0.2, 
                 gae_lambda=0.95):
        self.gamma = gamma
        self.eps_clip = eps_clip
        self.K_epochs = K_epochs
        self.gae_lambda = gae_lambda
        
        self.policy = ActorCritic(state_dim, action_dim)
        self.optimizer = optim.Adam(
            self.policy.parameters(), lr=lr, eps=1e-5)
        
        self.policy_old = ActorCritic(state_dim, action_dim)
        self.policy_old.load_state_dict(
            self.policy.state_dict())
    
    def compute_gae(self, rewards, values, dones, next_value):
        advantages = []
        gae = 0
        for t in reversed(range(len(rewards))):
            if t == len(rewards) - 1:
                next_val = next_value
            else:
                next_val = values[t + 1]
            
            delta = rewards[t] + self.gamma * next_val * \
                    (1 - dones[t]) - values[t]
            gae = delta + self.gamma * self.gae_lambda * \
                  (1 - dones[t]) * gae
            advantages.insert(0, gae)
        return advantages
\end{lstlisting}

\subsection*{附录B：训练参数配置}
\addcontentsline{toc}{subsection}{附录B：训练参数配置}

\begin{table}[H]
    \centering
    \caption{完整训练参数配置}
    \begin{tabular}{ll}
        \toprule
        \textbf{参数} & \textbf{值} \\
        \midrule
        环境 & CartPole-v1 \\
        最大训练轮数 & 2000 \\
        最大步数 & 500 \\
        更新步数 & 2048 \\
        打印频率 & 20 \\
        解决标准 & 平均得分$\geq$495 \\
        \bottomrule
    \end{tabular}
\end{table}

\end{document}
